\section{参考文献与阅读路线}

建议先用本讲义第 1--4 讲把舞台、主角、扭化与紧化这条主线走一遍，建立 $\tau,\Psi,t$ 与三种复结构之间的对应；再用第 5--7 讲把膜、线算子与 $D$-模的词典串起来。之后回到原始文献：先读 \cite{kw06} 的 \S1--\S5（约前三分之一），对照讲义逐个核对符号；再读 \S6--\S12，留意 't Hooft 算子的冒泡与分歧情形。物理背景可补读 \cite{w08} 的镜像与 Hitchin 视角，数学动机读 \cite{fr09} 的 Bourbaki 报告，分歧情形读 \cite{gw06}。最后可浏览 \cite{set13} 的博士论文与 \cite{bz04} 的课程笔记作为二次文献。

\begin{thebibliography}{9}
\bibitem{kw06} A.~Kapustin and E.~Witten, \emph{Electric--Magnetic Duality and the Geometric Langlands Program}, Commun.\ Number Theory Phys.\ 1 (2007) 1--236; arXiv:hep-th/0604151.
\bibitem{w08} E.~Witten, \emph{Mirror Symmetry, Hitchin's Equations, and Langlands Duality}, arXiv:0802.0999.
\bibitem{fr09} E.~Frenkel, \emph{Gauge Theory and Langlands Duality}, S\'eminaire Bourbaki 1010, Ast\'erisque 332 (2010) 369--403; arXiv:0906.2747.
\bibitem{gw06} S.~Gukov and E.~Witten, \emph{Gauge Theory, Ramification, and the Geometric Langlands Program}, arXiv:hep-th/0612073.
\bibitem{bz04} D.~Ben-Zvi, \emph{Between Electric--Magnetic Duality and the Langlands Program}, 课程笔记（arXiv 与作者主页可查）.
\bibitem{set13} N.~Setter, \emph{Topological Quantum Field Theory and the Geometric Langlands Program}, 博士论文.
\bibitem{ht03} N.~Hitchin, \emph{The Self-Duality Equations on a Riemann Surface}, Proc.\ London Math.\ Soc.\ 55 (1987) 59--126.
\end{thebibliography}
